\documentclass{article}
\usepackage[utf8]{inputenc}
\usepackage[T1]{fontenc}
\usepackage{geometry}
\usepackage{graphicx}
\usepackage{float}
\usepackage{hyperref}
\usepackage[english,indonesian]{babel}

\title{Optimasi Pemrosesan Template LaTeX Menggunakan Pendekatan Berbasis AI}
\author{Academic Support Agent}
\date{\today}

\begin{document}
\maketitle

\begin{abstract}

\end{abstract}

\chapter{PENDAHULUAN}


\chapter{TINJAUAN PUSTAKA}


% Original content (commented out):
% \documentclass[a4paper,12pt]{article}
    \include{01-bibliotecas}

\begin{document}

    \include{02-capa}
    \include{03-sumario} % é automático
    \include{04-texto} % o corpo do texto vai ficar aqui.
    
    %referências
    \titleformat*{\section}{\bf}
    \renewcommand{\refname}{REFERÊNCIAS}
    \addcontentsline{toc}{section}{REFERÊNCIAS}
    \printbibliography
    
\end{document}



\end{document}